\documentclass[11pt]{article}
\usepackage{fontspec}
\usepackage[french,english]{babel}
\usepackage{amsmath}
\usepackage{amssymb}
\usepackage{amsthm}
\usepackage{mathrsfs}
\usepackage{tikz-cd}
\usepackage{enumitem}

\newcommand{\printedpage}[1]{\marginpar{\footnotesize p.~#1}}
\newcommand{\msp}[1]{\marginpar{\footnotesize [#1]}}
\newcommand{\cardred}{\mathop{\mathrm{card}\,\textup{red}}\nolimits}
\newcommand{\Ind}{\mathrm{Ind}}
\newcommand{\circled}[1]{\tikz[baseline=(c.base)]\node[draw,circle,inner sep=1pt](c){#1};}

\begin{document}

\printedpage{44}
% head: A. Grothendieck : Dérivateurs, Ch. XVIII, Édition des Universités de Montpellier II et Paris VII
% foot: Version 17 décembre 2025

\msp{page 139}

\section*{4.6\quad Characterization of paratopos categories (proof of 4.3.6)}

It is time to return to the proof of Theorem 4.3.6 (p. 70-71), the one I had given being incomplete (and erroneous – I asserted that $(\mathrm{Filt}^{\pi})_{\pi\geqslant\pi_0}$ is a cardinal filtration, which is false if one assumes only $\pi^{\pi_0}=\pi$ for every $\pi\geqslant\pi_1$).

We had seen the equivalence of conditions a, a$'$, b, namely PTA. I am going to finish the proof by the logical scheme
\[
\text{PTA} \ \Longrightarrow\ \text{d)} \ \Longrightarrow\ \text{c)} \ \Longrightarrow\ \text{PTA}\ . \tag*{$(*)$}
\]

Let us prove PTA $\Longrightarrow$ d), i.e. that if $\mathcal{M}$ is an accessible pseudo-topos, then exists an infinite cardinal $\pi_0$, and a small category $C_0$ stable under $\varinjlim$ of cardinality $\leqslant\pi_0$, and an equivalence of categories % [sic] alors existe -> alors il existe, missing subject "il"
\[
\mathcal{M} \xrightarrow{\ \sim\ } \Ind'_{\pi_0}(C_0)\ . \tag*{(4.6.1)}
\]

More precisely, we have this:

\textbf{Proposition 4.6.1.} \textit{Let $\mathcal{M}$ be a paratopos, i.e. a $\mathfrak{U}$-category $\mathcal{M}$ stable under small $\varinjlim$ , admitting a full generating subcategory $S$ equivalent to a small category, consisting of accessible objects of $\mathcal{M}$. Let $\pi_0$ be a cardinal such that the objects of $S$ are $\pi_0$-accessible (i.e. $S\subset\mathcal{M}_{\pi_0}$), and let $C$ be the smallest strictly full subcategory of $\mathcal{M}$, containing $S$ and stable under $\varinjlim$ of cardinality $\leqslant\pi_0$. Then} % [sic] équivalent -> équivalente, agreement with the feminine "sous-catégorie"
\begin{itemize}[leftmargin=2.5em]
\item[\textit{1°)}] \textit{$C$ is equivalent to a small category.}

\msp{page 140}

\item[\textit{2°)}] \textit{$C=\mathcal{M}_{\pi_0}$, strictly full subcategory of $\mathcal{M}$ consisting of the $\pi_0$-accessible objects.}
\item[\textit{3°)}] \textit{The canonical functor}
\[
\Ind'_{\pi_0}(C) \xrightarrow[\sim]{\texttt{[}\varphi\texttt{]}} \mathcal{M} % ?? uncertain: font of the brackets around \varphi above the arrow -- typewriter (editorial) brackets or ordinary math brackets
\]
\textit{(defined up to unique isomorphism by the condition that its composite with the inclusion $C\hookrightarrow\Ind'_{\pi_0}(C)$ be isomorphic to the inclusion $C\longrightarrow\mathcal{M}$, and that $\varphi$ commute with small $\varinjlim$ , or again with filtering $\varinjlim$ large relative to $\pi_0$ – cf. 4.4.16, page 101) is an equivalence of categories} (\footnote[54]{thus we also have $\Ind'_{\pi_0}(C_0)\xrightarrow{\ \sim\ }\mathcal{M}$, where $C_0$ is a \textit{small} category equivalent to $C$.}).
\end{itemize}

\textsc{Proof}. 1°) We are going to construct $C$ by transfinite induction, setting
\[
\begin{aligned}
C_0 &\quad=\quad && \text{strictly full subcategory of $\mathcal{M}$ generated by $S$.}\\
C_{\alpha+1} &\quad=\quad && \text{strictly full subcategory of $\mathcal{M}$ generated by the $\varinjlim\limits_{I}$ in}\\
& && \text{$\mathcal{M}$ of objects of $C_\alpha$, with $I$ a category such that $\operatorname{card}\operatorname{Fl}I\leqslant\pi_0$.}\\
C_\alpha &\quad=\quad && \textstyle\bigcup_{\alpha'<\alpha}C_{\alpha'}\ \text{ if $\alpha$ is a limit ordinal.}
\end{aligned}
\]

This uniquely defines the $C_\alpha$ (for every small cardinal $\alpha$). By the definition of $C$, we see by transfinite induction that the $C_\alpha$ are contained in $C$. Let $\alpha_0$
\printedpage{45}
% head: A. Grothendieck : Dérivateurs, Ch. XVIII, Édition des Universités de Montpellier II et Paris VII
% foot: Version 17 décembre 2025
\msp{page 141}
% ?? uncertain: a hand-drawn mark (shaped roughly like a "2" or a bracket) stands in the left margin beside lines 3--5 of this page -- marginal annotation reproduced in the image, not typeset text
be the first ordinal such that $\operatorname{card}\alpha_0>\pi_0$, so the set $I_{\alpha_0}$ of cardinals $\alpha<\alpha_0$ is large relative to $\pi_0$. I say that we have $C_{\alpha_0}=C$. For this, it suffices to see that $C_{\alpha_0}$ is stable under $\varinjlim$ of cardinality $\leqslant\pi$. But if $I$ is a category such that $\operatorname{card}\operatorname{Fl}I\leqslant\pi_0$, and $\varphi : I\longrightarrow C_{\alpha_0}=\bigcup_{\alpha<\alpha_0}C_\alpha$ a functor, then the definition of $\alpha_0$ implies that $\varphi$ factors as $\varphi_\alpha : I\longrightarrow C_\alpha$, with $\alpha<\alpha_0$ (\footnote[55]{A bit short. One must invoke the fact that $I_{\alpha_0}$ is large relative to $\pi_0$.}). Thus $\varinjlim\varphi\in\operatorname{Ob}C_{\alpha+1}\subset\operatorname{Ob}C_{\alpha_0}$, which shows the required stability. % [sic] de cardinal \leqslant\pi -> de cardinal \leqslant\pi_0, the stability to be proved concerns cardinal \leqslant\pi_0

To see that $C$ is equivalent to a small category, it suffices to see this for the $C_\alpha$, proceeding by transfinite induction. This follows from the following lemma, immediate:

\textbf{Lemma 4.6.1.1.}
\begin{itemize}[leftmargin=2.5em]
\item[\textit{1°)}] \textit{The direct limit $C$ of a filtering direct system $(C_i)_{i\in I}$ of categories $C_i$ equivalent to small categories is equivalent to a small category. If $\cardred(C_i)$ $\leqslant\pi$ for every $i$, and $\operatorname{card}I\leqslant\pi$, then $\cardred(C)\leqslant\pi$ ($\pi$ being an} % [sic] système inductive -> système inductif, agreement with the masculine "système"

\msp{page 142}

\textit{infinite cardinal). Here, for any category $\mathcal{A}$, I put $\operatorname{card}\mathcal{A}=\operatorname{card}\operatorname{Fl}\mathcal{A}$, $\cardred(\mathcal{A})$ $=\operatorname{card}(\mathcal{A}_0)$, where $\mathcal{A}_0$ is a reduced category equivalent to $\mathcal{A}$.}
\item[\textit{2°)}] \textit{Let $C_0\subset\mathcal{M}$ be a full subcategory of a category $\mathcal{M}$, $\pi_0$ an infinite cardinal, and let $C'_0$ be the strictly full subcategory of $\mathcal{M}$ generated by the direct limits in $\mathcal{M}$ (representable in $\mathcal{M}$) of functors $I\longrightarrow C_0$, with $I$ a category such that $\operatorname{card}I\leqslant\pi_0$. If $C_0$ is equivalent to a small category, the same holds for $C'_0$. More precisely, if $\cardred(C_0)\leqslant\pi_1$, $\pi_1$ an infinite cardinal, then $\cardred(C'_0)\leqslant$ $\pi_1^{\pi_0}$} (\footnote[56]{This bound is not proved and probably false if one does not make somewhat stricter assumptions.}). % [sic] n'est par prouvée -> n'est pas prouvée, misprint in footnote 56
\end{itemize}

\textbf{Corollary 4.6.1.2.} \textit{$\mathcal{M}$ and $C_0$ being as in 2°) above, let $C$ be the smallest strictly full subcategory of $\mathcal{M}$ containing $C_0$, and stable under $\varinjlim$ of cardinality $\leqslant\pi_0$ (those which are representable in $\mathcal{M}$). If $C_0$ is equivalent to a small category, the same holds for $C$. More precisely, if $\cardred(C_0)\leqslant\pi_1$, where $\pi_1$ is an infinite cardinal, then}

\msp{page 143}

\textit{we have $\cardred(C)\leqslant\pi_1^{\pi_0}$} (\footnote[57]{bound not proved, probably false, cf. annotation above.}).

3°) is proved by the preceding transfinite proof, we prove by transfinite induction that
\[
\cardred C_\alpha \ \leqslant\ \pi_1^{\pi_0}
\]
for every $\alpha$. This is OK for $\alpha=0$ by hypothesis (since $\pi_1\leqslant\pi_1^{\pi_0}$). If it is true for $C_\alpha$, then 2°) ensures that
\[
\cardred C_\alpha \ \leqslant\ (\pi_1^{\pi_0})^{\pi_0} \ =\ \pi_1^{\pi_0\cdot\pi_0} \ =\ \pi_1^{\pi_0}\ , % [sic] C_\alpha -> C_{\alpha+1}, this is the successor step of the induction
\]
finally, if $\alpha$ is a limit ordinal $<\alpha_0$, we will have $C_\alpha=\bigcup_{\alpha'<\alpha}C_{\alpha'}$, hence
\[
\cardred C_\alpha \ \longrightarrow\ \pi_0\cdot\pi_1^{\pi_0} \ =\ \pi_1^{\pi_0} % [sic] \longrightarrow -> \leqslant, an inequality is intended
\]
(since $\pi_1^{\pi_0}\geqslant\pi_0$), by virtue of 1° above, and taking into account the fact that $\operatorname{card}I_\alpha\leqslant\pi_0$. Let $\pi'_0$ be the successor cardinal of $\pi_0$, then the same bound gives us
\[
\cardred C_{\alpha_0} \ \leqslant\ \pi'_0\pi_1^{\pi_0}\ .
\]
\printedpage{46}
% head: A. Grothendieck : Dérivateurs, Ch. XVIII, Édition des Universités de Montpellier II et Paris VII
% foot: Version 17 décembre 2025
But we have $\pi_1^{\pi_0}\geqslant2^{\pi_0}>\pi_0$, hence $\pi_1^{\pi_0}\geqslant\pi'_0$, hence $\pi_1^{\pi_0}\geqslant\pi'_0$, hence $\pi'_0\pi_1^{\pi_0}=\pi_1^{\pi_0}$, whence
\[
\cardred C_{\alpha_0} \ \leqslant\ \pi_1^{\pi_0}\ ,
\]
which is the announced bound.

\msp{page 144}

Let us prove 4.6.1 2°, i.e. $C=\mathcal{M}_{\pi_0}$. We have $C\subset\mathcal{M}_{\pi_0}$, since by hypothesis $S\subset\mathcal{M}_{\pi_0}$, and we know that $\mathcal{M}_{\pi_0}$ is stable under $\varinjlim$ of cardinality $\leqslant\pi_0$. On the other hand, we have $\mathcal{M}_{\pi_0}\subset C$, by virtue of 4.5.19 a $\Longrightarrow$ c applied to $C$, which is strictly generating in $\mathcal{M}$, $S$ being so, and $C\supset S$, and stable under $\varinjlim$ of cardinality $\leqslant\pi_0$ by construction.

\textit{Let us finally prove 3°}. This is equivalent to the following:

\textbf{Lemma 4.6.1.3.} \textit{Let $\mathcal{M}$ be a $\mathfrak{U}$-category stable under} small \textit{$\varinjlim$ , $\pi_0$ a cardinal, $\mathcal{M}_0$ a subcategory of $\mathcal{M}$} strictly generating, \textit{stable under $\varinjlim$ of cardinality $\leqslant\pi_0$, and consisting of $\pi_0$-accessible objects. (By virtue of 4.6.1 2° already proved, applied to $S=\mathcal{M}_0$, if $\mathcal{M}_0$ is strictly full, we therefore have $\mathcal{M}_0=\mathcal{M}_{\pi_0}$.) Then the canonical functor}
\[
\Ind'_{\pi_0}(\mathcal{M}_0) \ \longrightarrow\ \mathcal{M}
\]
\textit{is an equivalence of categories.}

The fact that this factor is fully faithful follows at once from the hypothesis that the objects of $\mathcal{M}_0$ are $\pi_0$-accessible: % [sic] facteur -> foncteur, misprint
\[
\begin{aligned}
\operatorname{Hom}_{\mathcal{M}}(\varinjlim_{i\in I}\underbrace{X_i}_{\in\mathcal{M}_0},\ \varinjlim_{j\in J}\underbrace{Y_j}_{\in\mathcal{M}_0}\ ) &\ \simeq\ \varprojlim_{i\in I}\ \underbrace{\operatorname{Hom}_{\mathcal{M}}(X_i,\ \varinjlim_{j\in J}Y_j)}_{\varinjlim\limits_{j\in J}\operatorname{Hom}(X_i,Y_j)}\\
&\ \simeq\ \underline{\operatorname{Hom}}_{\Ind_{\pi_0}(\mathcal{M})}((X_i),(Y_j))\ .
\end{aligned}
\]

\msp{page 145}

It remains to see that this functor is essentially surjective. Since $\mathcal{M}_0$ is strictly generating, it follows that every $x\in\mathcal{M}$ is written
\[
x \ =\ \varinjlim_{i\in I}x_i\ ,
\]
where $I=\mathcal{M}_0/x$, and the $x_i\in\mathcal{M}_0$. The hypothesis that $\mathcal{M}_0$ is stable under $\varinjlim$ of cardinality $\leqslant\pi_0$ implies the same property for $I=\mathcal{M}_0/x$, so $I$ is large relative to $\pi_0$. Thus $(x_i)_{i\in I}$ defines an object of $\Ind_{\pi_0}\mathcal{M}_0$, whose image in $\mathcal{M}$ is isomorphic to $x$, q.e.d. (Compare SGA 4 I 9.18.)

\textbf{Corollary 4.6.2.} \textit{Under the conditions of 4.6.1, let us take for $\pi_0$} the smallest \textit{cardinal such that the objects of $S$ are $\pi_0$-accessible. For every $\pi\geqslant\pi_0$, let $\mathcal{M}_\pi$ be the strictly full subcategory of $\mathcal{M}$ consisting of the $\pi$-accessible objects, which is also (by 4.6.1 2°) the closure of $C$ in $\mathcal{M}$ stable under $\varinjlim$ of cardinality $\leqslant\pi$.}
\begin{itemize}[leftmargin=2.5em]
\item[\textit{1°)}] \textit{Let $\pi_1=\cardred S$. We then have}
\[
\cardred\mathcal{M}_\pi \ \leqslant\ \pi_1^{\pi} \qquad (=2^{\pi}\ \textit{if}\ \pi\geqslant\pi_0). \tag*{(4.6.2.1)}
\]
\printedpage{47}
% head: A. Grothendieck : Dérivateurs, Ch. XVIII, Édition des Universités de Montpellier II et Paris VII
% foot: Version 17 décembre 2025
\msp{page 146}
\item[\textit{2°)}] \textit{$(\mathcal{M}_\pi)_{\pi\geqslant\pi_1}$ is a cardinal filtration of $\mathcal{M}$, where $\pi_1=\sup(\pi_0,\operatorname{card}S)$.}
\end{itemize}

\textsc{Proof}. The bound (4.6.2.1) is contained in 4.6.1.2, where $\pi_0$ has become $\pi$. If $\pi\geqslant\pi_1$, we have
\[
\pi_1^{\pi} \ =\ 2^{\pi} \tag*{$(*)$}
\]
for we have $\pi_1\geqslant2$, whence $\pi_1^{\pi}\geqslant2^{\pi}$, on the other hand, we have $\pi_1\leqslant2^{\pi_1}$, whence $\pi_1^{\pi}\leqslant(2^{\pi_1})^{\pi}=2^{\pi_1\pi}=2^{\pi}$, since $\pi_1\pi=\pi$ for $\pi\geqslant\pi_1$; whence $(*)$.

Assertion 2° is included for the record; by virtue of $\mathcal{M}\simeq\Ind'_{\pi_0}\mathcal{M}_{\pi_0}$, we may assume that we have $\mathcal{M}=\Ind'_{\pi_0}(\mathcal{M}_{\pi_0})$, and we can then apply 4.5.9. One would moreover like to say that $(\mathcal{M}_\pi)_{\pi\geqslant\pi_0}$ is a cardinal filtration, without restricting oneself to $\pi\geqslant\pi_1$. This is true except for the bound on cardinals in condition c) of the definition of cardinal filtrations (SGA 4 I 9.12):
\[
\underbrace{x}_{\in\operatorname{Ob}\mathcal{M}_{\pi'}} \ =\ \varinjlim_{i\in I}\underbrace{x_i}_{\in\operatorname{Ob}\mathcal{M}_{\pi}}\ , \qquad\qquad \text{$I$ large relative to $\pi$ and $\operatorname{card}I\leqslant\pi'^{\pi}$}
\]
for $\pi'\geqslant\pi$. One must assume, for this bound, that $\pi\geqslant\pi_1$.

\msp{page 147}

We can now finish the proof of 4.3.6, by proving the last two implications $(*)$ (p. 139). We have d $\Longrightarrow$ c, i.e. if $\mathcal{M}\simeq\Ind'_{\pi_0}(C)$, $C$ stable under direct limits of cardinality $\leqslant\pi_0$, and $C$ equivalent to a small category, then $\mathcal{M}$ is stable under small $\varinjlim$ (4.4.13, p. 95-96), its elements are accessible, and it admits a cardinal filtration (cf. 4.5.9). But one must check all the same that the elements of $C$ are $\pi_0$-accessible. This is contained in SGA 4 I 9.18, which asserts that the full faithfulness of the functor $\Ind_{\pi_0}(C)\longrightarrow\mathcal{M}$ \textit{is equivalent} to the fact that the elements of $C$ are $\pi_0$-accessible in $\mathcal{M}$. The proof in loc. cit. was not clear, the ``it suffices'' was seen bottom of page 144. The ``it is necessary'' does not seem clear to me, with the generality given in loc. cit. – and now I see that it is even false, as one sees by taking for $C$ the unit category $e$, so $\Ind'_{\pi_0}C\simeq e$, and the assertion affirming that if $e$ is fully embedded in $\mathcal{M}$ (thus by a \textit{rigid} object of $\mathcal{M}$, i.e. an object $a$ such that $\operatorname{End}a=\{\mathrm{id}_a\}$), then $a$ is $\aleph_0$-accessible. This is false in the category of rings, where $\mathbf{R}$ is a rigid object not $\aleph_0$-accessible!

\msp{page 148}

In the case we are concerned with, $\mathcal{M}=\Ind'_{\pi_0}C$, one must clearly use the explicit calculation of the $\varinjlim$ large relative to $\pi_0$ in $\Ind'_{\pi_0}C\subset C^{\wedge}$, they are indeed calculated ``argument by argument'', cf. 4.4.15.

This completes the proof of c $\Longrightarrow$ d in $(*)$, page 139. Finally the implication c $\Longrightarrow$ PTA is tautological, the existence of a cardinal filtration implying that of a small strictly generating subcategory. This completes the proof of 4.3.6.

I am going to give again a slightly more precise formulation of 4.3.6., specifying the cardinals, and adding to it an additional condition, inspired by SGA 4 I 9.11 (revised and corrected, cf. an annotation in the copy \texttt{[of Grothendieck]} of SGA 4 I (\footnote[58]{Ed.~note.\ \ This annotation says ``cf. simpler result above, if $C$ is \textit{strictly} generating'' and
% footnote 58 continues at the foot of printed page 48
the result in question, also in a marginal annotation in Grothendieck's copy, is:\par
\textbf{Corollary.} \textit{Let $F$ admit} \texttt{[a]} \textit{subcategory generating by strict epimorphisms, consisting of accessible objects, then} every \textit{object of $F$ is accessible (use 7.2. (i) b)).}})).

\printedpage{48}
% head: A. Grothendieck : Dérivateurs, Ch. XVIII, Édition des Universités de Montpellier II et Paris VII
% foot: Version 17 décembre 2025

\textbf{Theorem 4.6.3.} \textit{Let $\mathcal{M}$ be a $\mathfrak{U}$-category.}
\begin{itemize}[leftmargin=2.5em]
\item[\circled{1\textsuperscript{\textit{o}}}] \textit{The following conditions are equivalent for an infinite cardinal $\pi_0$:}
\begin{itemize}[leftmargin=3em]
\item[$a_{\pi_0}$\textit{)}] \textit{$\mathcal{M}$ is a $\mathfrak{U}$-$\pi_0$-paratopos category, i.e. $\mathcal{M}$ is stable under small $\varinjlim$ , and contains a small full} \texttt{[strictly]} \textit{generating subcategory $S$, consisting of $\pi_0$-accessible objects of $\mathcal{M}$.}

\msp{page 149}

\item[$b_{\pi_0}$\textit{)}] \textit{There exist a small category $S$, and a fully faithful and $\pi_0$-accessible functor $\mathcal{M}\xrightarrow{\ i\ }S^{\wedge}$, admitting a left adjoint.}
\item[$c_{\pi_0}$\textit{)}] \textit{There exists a small category $C$, stable under $\varinjlim$ of cardinality $\leqslant\pi_0$, such that $\mathcal{M}\simeq\Ind'_{\pi_0}(C)$.}
\item[$c'_{\pi_0}$\textit{)}] \textit{As c), with $C$ only equivalent to a small category.}
\end{itemize}
\item[\textit{2°)}] \textit{These conditions imply that $\mathcal{M}$ is stable under small $\varinjlim$ and small $\varprojlim$ , and the condition:}
\begin{itemize}[leftmargin=3em]
\item[$e_{\pi_0}$\textit{)}] \textit{For every category $I$ large relative to $\pi_0$, the functor}
\[
L \ =\ \varinjlim_{I}\ :\ \underline{\operatorname{Hom}}(I,\mathcal{M}) \ \longrightarrow\ \mathcal{M}
\]
\textit{is left exact (hence exact).}
\end{itemize}
\textit{This condition implies in particular the following:}
\begin{itemize}[leftmargin=3em]
\item[$\alpha(\pi_0)$] \textit{On the category of double arrows of $\mathcal{M}$, the functor} Ker \textit{(with values in $\mathcal{M}$) is $\pi_0$-accessible.}
\item[$\beta(\pi_0)$] \textit{For every arrow $y\leqslant x$ in $\mathcal{M}$, the base change functor $\mathcal{M}/x\longrightarrow\mathcal{M}/y$ is $\pi_0$-accessible.}
\end{itemize}
(\footnotemark[59]). % the text of footnote 59 is not printed on pages 46--50 of the PDF

\msp{page 150}

\item[\textit{3°)}] \textit{The following conditions are equivalent.}
\begin{itemize}[leftmargin=3em]
\item[\textit{a)}] \textit{$\mathcal{M}$ is a $\mathfrak{U}$-paratopos category, i.e. is stable under small $\varinjlim$ , and contains a small full subcategory $S$ generating by strict epimorphisms, and consisting of accessible objects.} % [sic] épimorphismes strictes -> épimorphismes stricts, agreement with the masculine "épimorphismes"
\item[\textit{a$'$)}] \textit{As a), but requiring that $\mathcal{M}$ be weakly accessible, i.e. that} every \textit{object of $\mathcal{M}$ be accessible.}
\item[\textit{b)}] \textit{There exist a small category $S$ and a fully faithful and accessible functor $\mathcal{M}\hookrightarrow S^{\wedge}$ admitting a left adjoint.}
\item[\circled{c}] \textit{There exist an infinite cardinal $\pi_0$, a small category $S$, and an equivalence $\mathcal{M}\simeq$ $\Ind'_{\pi_0}(S)$.}
\item[\textit{c$'$)}] \textit{As c, assuming only $S$ equivalent to a small category.}
\item[\textit{d)}] \textit{$\mathcal{M}$ is stable under small $\varinjlim$ , is weakly accessible, i.e. all its objects are accessible, and admits a cardinal filtration.}
\end{itemize}
\end{itemize}

\end{document}
